\emph{Within the value family}, adding a horizon does not help. The masked
Q-learner is the bandit with one change---a bootstrapped, discounted target
(Eq.~\ref{eq:qtarget}). With $\gamma=0.5$ it is close to the
bandit (mean difference \HQPointFiveVs, better on \HQPointFiveRootsBetter{} of
\HQPointFiveRoots{} roots); with $\gamma=0.9$ it is worse (\HQPointNine{} vs
\HQZero; mean difference \HQPointNineVs, better than the bandit on
\HQPointNineRootsBetter{} of \HQPointNineRoots{} roots), mainly through lower
network utility rather than higher churn (\HQPointNineRev{} reversals per
episode); with $\gamma=0.99$ it falls to \HQPointNineNine{} (mean difference
\HQPointNineNineVs, better on \HQPointNineNineRootsBetter{} of
\HQPointNineNineRoots{} roots). In both families the return decreases as the
discount grows beyond $0.5$. This matches the sequentiality diagnostic: the part of the future
that a learner could exploit from its observation is small, so a bootstrapped
target adds estimation error without adding much signal---the effect of a
shorter planning horizon under limited data~\citep{jiang2015horizon,amit2020discount}.

\emph{Shaping.} The training reward contains a potential-based shaping term
that is policy-invariant only for $\gamma=0.995$, so it could distort the
comparison across discounts. Removing it on two roots (42, 314159; evaluated
on the primary reward) leaves the ordering intact.
\ifcsname NoShapePPOVsBandit\endcsname Relative to the shaping-free bandit,
shaping-free PPO with $\gamma=0$ differs by \NoShapePPOZeroVsBanditPerRoot{},
the Q-learner with $\gamma=0.9$ by \NoShapeQVsBanditPerRoot{} and PPO with
$\gamma=0.995$ by \NoShapePPOVsBanditPerRoot{} (roots 42 / 314159). The
differences that survive learner noise (\S\ref{sec:rq3}) are the large ones,
and they point the same way with and without shaping.\fi
